% Q13_gold_answers.tex — GRADER CALIBRATION FIXTURE (correct, paraphrased).
% A drop-in replacement for paper/Q13_answers.tex (spec §4.7, §6.12): every answer is
% correct, but the route, the wording and the examples are deliberately far from the
% reference, so the grader must recognise correctness by substance (the rubric's
% probes), not by matching text. B1 uses pairs of its own for every missing property
% (none of the reference's) and argues (c) from what each side must prove; B2 words
% every field its own way.
%
% EXPECTED[B1][correctness][4]: correct; Probe A on each candidate agrees with every yes and no of (a), and Probe A' passes on every pair of (b): Q1 at (0, 3), Q2 at (6, 9), Q3 at (7, 0), Q5 at (-4, -4) for unsoundness and at (-4, 0) for incompleteness (captured 2026-10-11, Dafny 4.11.0: each `PASSED`; Probe A verified Complete for Q1-Q4 and Sound for Q4). No finding.
% EXPECTED[B2][card-fidelity][4]: every field agrees with the pinned range (lines 2005-2031); Budget and Variance correctly say what the range does not give. No finding on any field.
% Expected per-problem grades: B1 = 4, B2 = 4 (3 is tolerated only with a reading note and no flaw).
% A calibration run FAILS if any problem scores <= 2 or if any finding claims a false fact,
% a wrong pair, a wrong instrument outcome or a contradiction with the range.
\documentclass[11pt]{article}
\input{preamble}

\title{Q13 --- Practice II-A \\ \large Written answers}
\author{learner}
\date{}

\begin{document}
\maketitle
\sloppy

% ---------------------------------------------------------------------------
\section*{Problem B1 (concept)}\label{ans:B1}
\begin{answer}
Every true output of Ramp is non-negative and at least $x$, and is $x$ or $0$;
so the first three candidates all let every true output through.

$Q_1$: sound? no \quad complete? yes \quad pairs: at $x = 0$ it accepts
$y = 3$, while Ramp gives $0$. \\
$Q_2$: sound? no \quad complete? yes \quad pairs: at $x = 6$ it accepts
$y = 9$ ($9 \ge 0$, $9 \ge 6$), while Ramp gives $6$. \\
$Q_3$: sound? no \quad complete? yes \quad pairs: at $x = 7$ the disjunct
$y = 0$ accepts $0$, while Ramp gives $7$. \\
$Q_4$: sound? yes \quad complete? yes \quad pairs: none; for each sign of $x$
it names exactly Ramp's value. \\
$Q_5$: sound? no \quad complete? no \quad pairs: at $x = -4$ it accepts
$y = -4$ (Ramp gives $0$), and it rejects Ramp's own $(-4, 0)$. \\
\textbf{(c)} A caller of Ramp can only use what the postcondition promises, so
it is best served by the sound and complete $Q_4$; the body that must
establish the postcondition has least to prove under the loosest one, $Q_1$.
Strength is a benefit on the calling side and a cost on the proving side.
\end{answer}

% ---------------------------------------------------------------------------
\section*{Problem B2 (paper card)}\label{ans:B2}
\begin{answer}
\textbf{Source:} cheatsheet entry 05.5, lines 2005--2031 (the pinned range).

\textbf{Task shape:} 189 Lean 4 problems, checked by hand, where a system must
write the program, its pre- and postcondition, and a proof that they hold; two
splits, one translated from Dafny (108) and one from course submissions (81).

\textbf{Oracle:} The Lean kernel for code and proofs. For a generated
specification, soundness and completeness against the task's reference
specification are tried by a proof from the model, then by deciding test
instances, then by property-based tests; anything still open is ``unknown''.

\textbf{What is trusted:} Lean itself, one reference specification per task,
and the tests that come with it (positive and negative, covering every line).

\textbf{Cheat surface:} Writing a specification that mirrors the code would
make a proof easy, so end-to-end proofs must target the ground-truth
specification; the test-based spec metric leaves up to roughly 30\% of hard
postconditions unknown.

\textbf{Metric:} Pass@1 for code, specification, proof and the whole chain:
o3 reaches 72.6, 52.3, 4.9 and 3.2. Fewer than 4\% of specifications are
certified by proof, versus 41.6 to 51.0\% by testing.

\textbf{Budget:} The range gives a token budget only for the refinement
comparison with COPRA (o4-mini, 350k tokens) and 64 refinement rounds for
DeepSeek Prover V2; total compute is not stated.

\textbf{Variance:} The range reports no repeated runs or spreads; I would need
seeds or confidence intervals.

\textbf{Artifact:} The VERINA benchmark, both splits.
\end{answer}

\end{document}
